\chapter{Acoustic Neuroma}

\section*{Definition and Overview}
An acoustic neuroma (also called a vestibular schwannoma) is a slow-growing, usually non-cancerous (benign) tumour that develops on the nerve responsible for balance and hearing as it travels from the inner ear to the brainstem. It arises from the Schwann cells that form the insulating layer around the nerve fibres.

Despite its traditional name, the tumour usually arises from the vestibular (balance) nerve rather than the acoustic (hearing) nerve, which lies alongside it. As it grows it can compress the hearing and balance nerves and, if it becomes large, the brainstem, causing significant hearing and balance problems and, rarely, dangerous pressure inside the skull.

\section*{Epidemiology}
Acoustic neuromas are uncommon but are being diagnosed more often, partly because of the wider use of MRI scanning. They affect roughly 1 to 2 people per 100,000 each year, are most commonly diagnosed between the ages of 40 and 60, and affect men and women about equally.

The vast majority are unilateral (affecting one side only) and arise for no identifiable reason. Bilateral tumours are rare and are a feature of the inherited condition neurofibromatosis type 2 (NF2), which usually presents at a younger age.

\section*{Aetiology and Pathophysiology}
In most cases the cause is unknown, and the tumour develops from a spontaneous genetic change in a single Schwann cell. There is no established link to everyday factors such as mobile phone use, and the current evidence does not support such a connection.

In the inherited form, a faulty \textit{NF2} gene is passed down in families, allowing Schwann cells to multiply out of control, often on both sides. As the tumour grows within the narrow bony canal carrying the nerve, it first compresses the balance and hearing nerves. Larger tumours can extend further and press on the brainstem and on the nerves that control facial movement and sensation, causing additional symptoms.

\section*{Clinical Features}
\begin{itemize}
    \item \textbf{Hearing loss:} progressive, gradual hearing loss on one side, the most common first symptom
    \item \textbf{Tinnitus:} ringing, buzzing, or hissing in one ear
    \item \textbf{Ear fullness:} a feeling of pressure or blockage in the affected ear
    \item \textbf{Balance disturbance:} dizziness, imbalance, or unsteadiness, though true spinning vertigo is less common
    \item \textbf{Speech understanding:} difficulty understanding speech, even when sounds can still be heard
    \item \textbf{Large-tumour features:} facial numbness or weakness, headache, and unsteadiness of walking
\end{itemize}

\section*{Differential Diagnosis}
Unilateral hearing loss, tinnitus, and dizziness have many other possible causes.

\begin{itemize}
    \item Meniere's disease (episodes of vertigo, hearing loss, and tinnitus)
    \item Vestibular neuronitis or labyrinthitis (inflammation of the balance structures)
    \item Middle-ear infections and fluid
    \item Noise-induced or age-related hearing loss
    \item Sudden sensorineural hearing loss from other causes
    \item Other tumours of the same region, such as meningiomas
    \item Multiple sclerosis affecting the brainstem
\end{itemize}

\section*{When to Seek Medical Care}
See a doctor promptly for any new hearing loss, ringing, or dizziness affecting one ear, especially if it is progressive or persistent. Sudden hearing loss needs early assessment. Anyone with a known acoustic neuroma should report new or worsening symptoms, especially facial weakness, headache, or balance problems.

\textbf{Call 000 (or local emergency services) for:}
\begin{itemize}
    \item Sudden severe headache with vomiting
    \item Sudden loss of vision, double vision, or confusion
    \item Sudden weakness or numbness of the face, arm, or leg
    \item Difficulty walking or standing that comes on suddenly
\end{itemize}

\section*{Diagnosis and Investigations}
\begin{itemize}
    \item \textbf{Hearing tests (audiometry):} measure hearing in each ear and how well speech is understood
    \item \textbf{MRI scan with contrast:} the definitive test, which shows the tumour and its exact size and position
    \item \textbf{CT scan:} may be used where MRI is not possible or for planning surgery
    \item \textbf{Balance testing and facial nerve assessment:} important for larger tumours
    \item \textbf{Monitoring:} small, slow-growing tumours are often followed with repeated MRI scans over time
\end{itemize}

\section*{Management}
Treatment depends on the size of the tumour, its growth rate, the person's age and general health, the degree of hearing loss, and personal preference.

\begin{itemize}
    \item \textbf{Observation (watch and wait):} small, slow-growing tumours in older people or those with minimal symptoms may simply be monitored with regular MRI scans
    \item \textbf{Stereotactic radiosurgery:} precisely focused radiation (often called gamma knife) that aims to stop the tumour growing without an open operation
    \item \textbf{Surgery:} removal of the tumour, usually recommended for larger tumours or those causing pressure; hearing on that side is often lost
    \item \textbf{Hearing rehabilitation:} hearing aids, and consideration of special devices or implants for people with total hearing loss
    \item \textbf{Balance and facial rehabilitation:} physiotherapy and specialist support for ongoing symptoms
    \item \textbf{Genetic advice:} for bilateral tumours or a family history, referral for genetic counselling for NF2
\end{itemize}

\section*{Complications}
\begin{itemize}
    \item Permanent hearing loss in the affected ear, which may occur even without treatment
    \item Persistent tinnitus
    \item Facial weakness or numbness, especially after surgery for larger tumours
    \item Risks of surgery, including infection, bleeding, fluid leakage from the ear, and damage to nearby nerves
    \item Balance problems and a lasting sense of unsteadiness
    \item Growth of an observed tumour, requiring a change in treatment
\end{itemize}

\section*{Prognosis and Outcomes}
The outlook is generally good, because acoustic neuromas grow slowly and are almost always benign. Many small tumours grow very little or not at all over many years, and observation is a safe and common option, especially in older people. When treatment is needed, both radiosurgery and surgery control growth in the large majority of cases.

Hearing on the affected side is often lost regardless of the approach, and there is a small risk of facial weakness with surgery. Regular MRI follow-up is important because treatment does not guarantee that the tumour will never grow again, but most people continue to lead full and active lives with support for hearing and balance.

\section*{Prevention and Self-Care}
\begin{itemize}
    \item Protect your hearing with ear protection in noisy environments and safe listening volumes
    \item Report any new unilateral hearing loss, ringing, or dizziness early
    \item Attend all scheduled follow-up MRI scans if you are under observation
    \item If you have hearing loss, use a hearing aid and seek support for communication
    \item Manage balance symptoms with physiotherapy and by avoiding sudden movements
    \item For families with NF2, seek genetic counselling before planning a pregnancy
\end{itemize}

\section*{Sources and Further Reading}
The following organisations provide reliable information on acoustic neuromas and vestibular schwannomas.

\begin{itemize}
    \item NHS. \textit{Information on acoustic neuroma, diagnosis, and treatment.} https://www.nhs.uk/conditions/
    \item Mayo Clinic. \textit{Guide to vestibular schwannoma and its management.} https://www.mayoclinic.org/
    \item MedlinePlus. \textit{Health information on acoustic neuroma and ear disorders.} https://medlineplus.gov/
    \item MSD Manual. \textit{Professional information on tumours of the ear and cranial nerves.} https://www.msdmanuals.com/
    \item NICE. \textit{Clinical guidance relevant to brain tumours and neurosurgery.} https://www.nice.org.uk/
\end{itemize}
